\section{Monte Carlo 用随机样本逼近期望}

一家电网公司想估计明年的期望停电损失。负载、天气和设备状态共同决定损失函数 \(h(X)\)，即使这些变量的联合分布 \(p(x)\) 已知，积分

\[
\mu=\E_p[h(X)]
=\int h(x)p(x)\,dx
\]

也可能因维度高和 \(h\) 非线性而无法解析求出。规则网格在每个维度取 \(m\) 个点需要 \(m^d\) 次计算；Monte Carlo 改用随机样本。

若

\[
X_1,\ldots,X_n\overset{\mathrm{iid}}{\sim}p,
\]

则估计

\[
\widehat\mu_n
=\frac1n\sum_{i=1}^nh(X_i).
\]

它把积分问题转成了平均问题。只要 \(\mathbb E|h(X)|<\infty\)，大数定律给出一致性；记 \(\sigma_h^2=\operatorname{Var}(h(X))\)，若 \(\sigma_h^2<\infty\)，则

\[
\operatorname{Var}(\widehat\mu_n)
=\frac{\sigma_h^2}{n},
\]

并由中心极限定理得到

\[
\sqrt n(\widehat\mu_n-\mu)
\overset d\longrightarrow
\mathcal N(0,\sigma_h^2).
\]

样本标准差 \(s_h\) 因此给出 Monte Carlo 标准误 \(s_h/\sqrt n\)。

\subsection{无偏不等于精确，计算误差也要进入报告}

\(\widehat\mu_n\) 无偏，只表示在无数次独立运行中平均正确，不表示某一次模拟已经足够接近。误差按 \(n^{-1/2}\) 缩小：若希望标准误减半，通常需要四倍样本；多打印几位小数并不会增加信息。

对电网损失，除了期望，还可能关心超过容量阈值的概率

\[
P(h(X)>c)
=\E[\mathbf 1\{h(X)>c\}].
\]

若事件非常罕见，普通模拟可能几乎看不到失败，估计值甚至为零。这不意味着风险为零，而是采样分布把计算资源浪费在大量普通情形上。

\subsection{方差缩减比盲目增样本更聪明}

importance sampling 从另一个分布 \(q\) 采样：

\[
\E_p[h(X)]
=\mathbb E_q
\left[
h(X)\frac{p(X)}{q(X)}
\right].
\]

把右端期望按定义展开即 \(\int h(x)\frac{p(x)}{q(x)}q(x)\,dx=\int h(x)p(x)\,dx\)，所以这只是同一个积分的改写，前提是约去 \(q(x)\) 合法。若 \(q\) 更常访问对积分贡献大的区域，方差可显著下降。但必须满足目标有贡献的地方 \(q(x)>0\)；否则某些区域永远采不到，估计发生不可修复的偏差。即便支持覆盖，极端 importance weights 也会导致巨大甚至无穷方差：权重越集中，有效样本量 \((\sum_iw_i)^2/\sum_iw_i^2\) 越小，表面样本量再大也主要由少数样本决定。

control variate 利用一个与 \(h(X)\) 相关、且期望已知的量 \(g(X)\)：

\[
\widehat\mu_{\mathrm{cv}}
=\frac1n\sum_i
\left[
h(X_i)-c(g(X_i)-\mathbb Eg)
\right].
\]

若选择适当 \(c\)，\(g\) 与 \(h\) 的共同波动被抵消。修正项期望为零，估计保持无偏；把方差写成 \(c\) 的函数，即 \(n^{-1}\bigl(\sigma_h^2-2c\operatorname{Cov}(h,g)+c^2\sigma_g^2\bigr)\)，这是关于 \(c\) 的二次函数，在 \(c^*=\operatorname{Cov}(h,g)/\sigma_g^2\) 处取最小，\(g\) 与 \(h\) 相关越强，方差下降越多。分层抽样保证不同区域都被覆盖，antithetic sampling 则利用负相关样本抵消波动。它们没有改变目标期望，而是在相同计算预算下减少方差。

\subsection{什么时候中心极限定理不能救场}

若 \(h(X)\) 重尾到方差无穷，常见的 \(\pm1.96\,s/\sqrt n\) 误差区间没有通常依据。importance weights 的重尾也可能让表面样本量很大、有效信息却由少数样本控制。此时应检查批次稳定性、最大权重、尾部诊断，并考虑重新设计提议分布或改变估计方法。

一个可靠模拟应记录随机种子、样本数、点估计和 Monte Carlo 标准误；用多批独立运行检查误差是否大致按 \(n^{-1/2}\) 缩放；在存在解析答案或低维数值积分的简化版本上先核对实现。模拟误差与模型不确定性也要分开：把 Monte Carlo 样本加到无限多，只会更精确地算出当前模型下的答案，不能修复错误的天气分布或损失函数。
